\section{From reciprocal collisions to local occupation queries}
\label{sec:queries}
This section supplies the observation reduction used by every adaptive
query. The random walk is computational: its exit time is not an
observed collision record or an apparatus stopping signal.

\begin{lemma}[Reciprocal balance and calibration]\label{lem:balance}
With $\chi=\one_{\mathcal O}$, outside null boundary sets,
\begin{equation}\label{eq:balance}
 B_a(q)-B_{-a}(q+a)=\chi(q+a)-\chi(q).
\end{equation}
Consequently the difference $g_\sigma$ of the two pooled nominal means
satisfies
\begin{equation}\label{eq:forcing}
 g_\sigma=(T-I)u_\sigma,\qquad
 Tv(x)=\tfrac14\sum_{a\in\{\pm te_1,\pm te_2\}}v(x+a).
\end{equation}
For either command, coupled errors of size
$r=\ell+\tau+t\alpha\le\sigma$ change its mean by at most $Cr/\sigma$.
The constant is uniform over the protected aperture, components and
nominal directions, including tangencies.
\end{lemma}
\begin{proof}
If both endpoints are free, a segment and its reverse either both hit
or both miss. If exactly one endpoint is solid, the direction starting
there gives zero and the other direction hits. If both are solid both
bits are zero. This proves \eqref{eq:balance}. Averaging in $q,a$ and
commuting translation with convolution proves \eqref{eq:forcing}.

For a convex body put $S_a(C)=C+[-a,0]$. The hit bit is
\[
 B_a(q)=\one_{\cup_C S_a(C)}(q)-\chi(q).
\]
The displacement error is at most $\tau+t\alpha$. Support inequalities
show that membership can change only if the nominal start lies in
an $r$-tube of $\partial C$ or $\partial S_a(C)$, including the position
error. The part of a convex boundary tube in a disk of radius $\sigma$
has area at most $C(\sigma+r)r$. One proof integrates boundaries of
outer parallel bodies and inner erosions by the coarea formula; their
portions in the disk have length at most $2\pi\sigma$, by monotonicity
of the perimeter of convex sets. Degenerate levels follow by a limit.
Only a bounded number of components contribute: each has a point in
a disk of radius $\sigma+t+r$, and such points from distinct components
are $d_0$-separated. Multiplication by the density bound
$\sigma^{-2}\|\kappa\|_\infty$ proves the assertion. The tube containment
is pointwise for every allowed error, so the coupling need not have
mean zero. An exceptional coupling event adds its probability.
\end{proof}

Put $D=D_0+2\sigma_0$, $H=(D+t)^2/t^2$ and
$L_* =\max(1,\lceil2H\rceil)$. The positive set of $u_\sigma$ lies in
the $\sigma$-enlarged components. They have diameter at most $D$ and
mutual distance greater than $t$.

\begin{lemma}[Stopped inverse in one aperture]\label{lem:stopped}
For the compass walk $X_n=x+a_1+\cdots+a_n$ and
$\tau_u=\inf\{n\ge0:u_\sigma(X_n)=0\}$ one has
\begin{equation}\label{eq:survival}
 \sup_x\E_x\tau_u\le H,\qquad
 \sup_x\Prob_x(\tau_u>n)\le2^{-\lfloor n/L_*\rfloor}.
\end{equation}
Let a bounded target region $U$ contain the complete protected bodies
and their search brackets. Choose a fixed bounded $W$ containing
$U+\overline B_{D+t+1}$. On $W$ use the killed transition
$T^Wv=T(\one_Wv)$ and set iterates to zero outside $W$. For
$V_0=0$, $V_{n+1}=\max(0,T^WV_n-g_\sigma)$,
\begin{equation}\label{eq:bellman-error}
 0\le u_\sigma-V_n\le2^{-\lfloor n/L_*\rfloor}\quad\hbox{on }U.
\end{equation}
Uniform forcing error $\xi$ and numerical update error $\zeta$ add
at most $n(\xi+\zeta)$ if approximate iterates are clipped to $[0,1]$.
\end{lemma}
\begin{proof}
Until its first zero, a path stays in one enlarged component, since
one step cannot reach another. For $S=\tau_u\wedge n$,
$|X_S-x|\le D+t$. The bounded stopping identity for the martingale
$|X_n-x|^2-nt^2$ gives $t^2\E S\le(D+t)^2$. Monotone convergence
proves the mean bound. Markov's inequality at $L_*$ and the Markov
property at successive blocks prove the geometric tail.

On $W$, $T^Wu_\sigma\le Tu_\sigma$, so $0\le V_n\le u_\sigma$
and the iterates vanish on its zero set. Starting in $U$, the complete
containing component and every first-exit location lie in $W$. Along
this stopped path the killed transition has no boundary discrepancy.
For $e_n=u_\sigma-V_n$, the recursion gives $e_{n+1}\le Te_n$ before
exit, and hence
\[
 e_n(x)\le\E_x[u_\sigma(X_n)\one_{\{\tau_u>n\}}].
\]
This proves \eqref{eq:bellman-error}. Maximum, averaging and clipping
are nonexpansive in supremum norm; induction proves the error bound.
No periodic wrap or renormalization of a boundary row is used.
\end{proof}

The stopped argument also compares two physical occupations with
unrelated zero sets. Indeed, for $w=u-\widetilde u$ and
$h=(T-I)w$, bounded stopping at $\tau_u\wedge n$ gives
\[
 w(x)=\E_xw(X_{\tau_u\wedge n})
                 -\E_x\sum_{j<\tau_u\wedge n}h(X_j).
\]
At exit the terminal value is $-\widetilde u\le0$. On survival it
is at most one. Letting $n$ increase and then interchanging $u$ and
$\widetilde u$ yields $\|u-\widetilde u\|_{U,\infty}
\le H\|(T-I)(u-\widetilde u)\|_{W,\infty}$. Thus exact reciprocal
functionals at all sufficiently small scales determine the local
occupation and, by approximate-identity convergence and regular
closedness, the physical set. This retained exact conclusion requires
neither analyticity nor a supplied zero set. The martingale and killed
Green mechanisms are classical; see \cite{LawlerLimic}.

\begin{proposition}[A finite local query from pooled bits]
\label{prop:query}
For any adaptively chosen $x\in U$ and $\sigma\le\sigma_0$, an
occupation estimate with error less than $1/8$ can be computed using
only a prior-bounded number of reciprocal command centers in $W$.
At conditional confidence $1-\varepsilon$ it uses at most
$C\log(C/\varepsilon)$ attempted preparations. A sufficient calibration
is $\ell+\tau+t\alpha\le c\sigma$. Thresholding the estimate at
$1/2$ gives the correct membership label outside the $\sigma$-tube
of the component boundaries. Its label inside that tube is unrestricted.
\end{proposition}
\begin{proof}
Fix $n_*=6L_*$. To evaluate $V_{n_*}(x)$, its dependency locations
are exactly among
\begin{equation}\label{eq:diamond}
 x+t(i,j),\qquad |i|+|j|\le n_*.
\end{equation}
Intersect these with $W$ and put zero at locations outside it. Evaluate
the recursion backwards through shrinking diamonds, initializing
$V_0=0$. Induction on the layer shows equality with the full killed
iterate at the target. There are $O(n_*^2)$ centers and $O(n_*^3)$
updates, both independent of $\sigma$. No fine background spatial grid,
interpolation of the transition, or growing aperture is involved.

At each center estimate each pooled mean to error at most
$1/(1024n_*)$, and impose the same bound on its calibration bias.
The forcing error is then at most $1/(256n_*)$. Reserve update error
at most $1/(256n_*)$. Lemma~\ref{lem:stopped} gives total error at most
$1/64+1/256+1/256<1/8$. Hoeffding's inequality and a union bound over
the fixed finite number of means give the preparation bound. The
constant includes $n_*$ and may be large. The proof holds conditional
on past data because the new attempts are fresh. Inside a component
at boundary distance greater than $\sigma$ the average is one;
outside all enlarged components it is zero. Thresholding therefore
has the stated property. A membership label is a computed conclusion
from these bits, not an extra response furnished by the sensor.
\end{proof}

The full launch region may be taken to contain
$W+\overline B_{t+\sigma_0+1}$, which includes reverse positions and
kernel supports. For at most $Q$ adaptive queries, allocating conditional
failure probability $\delta/Q$ to each gives joint correctness with
probability $1-\delta$. It suffices to know a deterministic upper bound
on $Q$; repeated centers may be resampled. No independence of the
adaptively selected locations themselves is asserted.
